% Small replacements for Online Appendix passages outside the M2 subsection list that
% restate the low cost, its floor condition, or the retired price-pool family.
% Each block carries its own header. Where the general-model rewrite also rewrites
% the passage, the integrator takes the substance and drops the block.
% Numbering note for every block that cites (A1)-(A3): these blocks keep the surviving
% conditions under their current main-text numbers (A2) and (A3). If the main text
% renumbers them after deleting (A1), the integrator renames them here in the same pass.

% REPLACES: online_appendix.tex lines 206-208 ("The challenger observes P but not X ..." to "... Paying the cost")
The challenger observes \(P\) but not \(X\), then its independent
preparation cost, equal to \(0\) with probability \(\rho\) and \(c_H\)
otherwise, where \(0<\rho<1\) and \(c_H>0\). Paying the cost

% REPLACES: online_appendix.tex lines 222-233 ("m=(1+e^{2/b})^{-1}, and M=1-m. ..." to "... (OA.1)")
\(m=(1+e^{2/b})^{-1}\), and \(M=1-m\). For the comparison
\(p<\ell<r_0<r_1<h\), the hypotheses of Proposition 2 are

\[
\begin{aligned}
&B_{r_0}(1/2)<c_H<B_{r_1}(M),\\
&\Delta_T(r_0)<k<(1-1/b)\rho m\Delta_T(r_1),
\end{aligned}
\tag{OA.1}\label{eq:oa-oa-1}
\]

% REPLACES: online_appendix.tex lines 297-299 ("The bidder's preparation rule is ..." to "... bounded below by \rho.")
The bidder's preparation rule is \(\mathbf1\{C\le B_r(\mu_P)\}\), with
entry at equality. The cost-averaged rule is
bounded below by \(\rho\) at every belief, because the zero-cost
realization prepares whenever gross profit is nonnegative, and gross
profit is positive at each feasible belief on this support. Conditional on \(\Theta\) and \(X\), the

% REPLACES: online_appendix.tex line 379 ("For any candidate equilibrium with the low-cost entry floor, ...")
For any candidate equilibrium, the

% REPLACES: online_appendix.tex lines 490-492 ("\textbf{Weak incumbent.} Proposition A.1 and (A1) give ..." to "... gives the residual upper bound.")
\textbf{Weak incumbent.} Proposition A.1 and the zero preparation cost give the entry floor
under every conditional mixed order. \hyperref[oa-a-inference]{Section
A.3} gives the residual upper bound.

% REPLACES: online_appendix.tex lines 511-515 ("Condition (A2) excludes high-cost ..." to "... preparation outcome.")
\(P=t_0+\rho[(t_H+t_L)/2-t_0]\). Condition (A2) excludes high-cost
preparation at the prior, and the zero-cost realization prepares there.
Competitive pricing is correct and the previous payoff calculation
validates zero. This constructs the unique trading and on-path
preparation outcome.

% REPLACES: online_appendix.tex lines 622-627 ("For Proposition A.4, define the maintained domain ..." to "... (OA.19)")
For Proposition A.4, define the maintained domain

\[
\mathcal D=\{r\in(\ell,h):B_r(1/2)<c_H\}.
\tag{OA.19}\label{eq:oa-oa-19}
\]

% REPLACES: online_appendix.tex line 663 ("the root residual, and the support and floor conditions ...")
% One phrase only; the sentence around it stands.
the root residual, and the support conditions, including \(B_r(m)>0\) at the node, rather than

% REPLACES: online_appendix.tex lines 683-687 ("For Proposition 2(iii), impose (A1) to (A3) ..." to "... strengths between them.")
For Proposition 2(iii), impose the conditions in \eqref{eq:oa-oa-1} at \(r_0,r_1\) and \(B_{r_2}(M)<c_H\) at \(r_2\); the global derivative bound at \(r_2\) follows from (A3) because \(\Delta_T\) increases in \(r\). Parts (i) and (ii) of Proposition 2 give the first
two unique outcomes; the bound and infeasibility condition give unique
full orders but entry \(\rho\) at \(r_2\). This proves a finite rise and
fall without a selection or uniqueness assertion for strengths between
them.

% REPLACES: online_appendix.tex lines 787-792 ("The quantity B_{r_1}(m) is positive, so choose c_L ..." to "... at every value ratio.")
The quantity \(B_{r_1}(m)\) is positive, so the zero cost lies below
\(c_H\): \(0<B_{r_1}(m)<B_{r_0}(1/2)<c_H\), using \(m<1/2\) and the monotonicity in
both arguments. All inequalities have strict slack and persist by
continuity. This is a mathematical nonemptiness construction, not a
statement about empirical effect size at every value ratio.

% REPLACES: online_appendix.tex lines 798-803 ("\protect\phantomsection ... Lemma OA.1 ..." to "... remain valid.")
\protect\phantomsection\label{result-lemma-oa-1}{\textbf{Lemma OA.1
(analytical).}} \emph{Replace the trading cost by
\(K(s)=ks+\gamma_qs^2/2\), \(\gamma_q\ge0\). Under (A2),
\(\Delta_T(r_0)<k\), and \(k+\gamma_q<(1-1/b)\rho m\Delta_T(r_1)\), the
unique trading outcomes and entry reversal in parts (i) and (ii) of
Proposition 2 remain valid.}

% REPLACES: online_appendix.tex lines 813-820 ("\protect\phantomsection ... Lemma OA.2 ..." to "... unique trading outcomes hold.")
% Integrator: merge with the general-prior notation of the main rewrite
% (m(\pi), M(\pi), \kappa(\pi)). The substance for the zero cost is the deleted
% first condition and the sentence after the conditions.
\protect\phantomsection\label{result-lemma-oa-2}{\textbf{Lemma OA.2
(analytical).}} \emph{For prior \(\pi\in(0,1)\), let
\(\mu_-=\operatorname{logistic}(\operatorname{logit}\pi-2/b)\),
\(\mu_+=\operatorname{logistic}(\operatorname{logit}\pi+2/b)\), and
\(\omega=\min\{\mu_-,1-\mu_+\}\). Suppose
\(B_{r_0}(\pi)<c_H<B_{r_1}(\mu_+)\) and
\(\Delta_T(r_0)<k<(1-1/b)\rho\omega\Delta_T(r_1)\). Then the benchmark
entry reversal and unique trading outcomes hold. No condition on the low cost
is needed: gross profit at \(\mu_-\) is positive on this support, so the
zero-cost realization prepares at every feasible belief.}

% REPLACES: online_appendix.tex lines 1827-1832 ("establish the stated result, I require ..." to "... at the same strength.")
% The pooling-existence margin later in the sentence is the zero-order margin and stays.
establish the stated result, I require all of the following: the model
support inequalities, zero-cost participation read as \(B_r(m)>0\) and high-cost prior
exclusion, strict threshold ordering throughout the bracket, opposite
enclosed endpoint signs in \eqref{eq:oa-oa-70}, positive \(\Gamma_H\) in
\eqref{eq:oa-oa-73}, and a positive pooling-existence margin at the same
strength.

% REPLACES: online_appendix.tex lines 2107-2113 ("\textbf{Inputs.} Use the benchmark declaration ..." to "... define separate economies.")
% Numbers: \NUM{atomless_low_upper} and the law name come from the numerics report.
% The text below states only the supports, so it holds for the uniform law or any
% replacement law on the same supports.
\textbf{Inputs.} Use the benchmark declaration at weak, strong, and
collapse strengths. Evaluate the atomic-cost model with Laplace noise
first. Then replace the noise law by logistic noise, retaining the scale
parameter, and replace cost atoms by the declared atomless components on
\([0,\NUM{atomless_low_upper}]\) and
\([c_H-\varepsilon_C,c_H+\varepsilon_C]\). Their mixture weights remain
\(\rho,1-\rho\). These replacements define separate economies.

% REPLACES: online_appendix.tex lines 2118-2134 ("B_r(m),B_r(1/2),B_r(M). Theorem margins are ..." to "... assigning an analytical equilibrium label.")
\(B_r(m),B_r(1/2),B_r(M)\). Theorem margins are

\[
\begin{aligned}
\zeta_L&=B_{r_1}(m),\qquad
\zeta_{H0}=c_H-B_{r_0}(1/2),\qquad
\zeta_{H1}=B_{r_1}(M)-c_H,\\
\zeta_0&=k-\Delta_T(r_0),\qquad
\zeta_1=(1-1/b)\rho m\Delta_T(r_1)-k.
\end{aligned}
\tag{OA.77}\label{eq:oa-oa-77}
\]

Here \(\zeta_L\) is gross profit at the lowest feasible belief. It is the margin
by which the zero-cost realization prepares at that belief, and a recorded
quantity rather than a hypothesis on a cost. The collapse node additionally requires \(c_H-B_{r_2}(M)>0\),
\(B_{r_2}(m)>0\), and its full-order margin. For atomless costs
replace the appropriate cost endpoints as in Proposition A.6: the low margin
then reads \(B_{r_1}(m)-\varepsilon_C\). Verify all
margins before assigning an analytical equilibrium label.

% REPLACES: online_appendix.tex lines 2294-2295 ("\textbf{Analytically classified outcomes.} First verify membership ..." to "... at each node. Pooling existence is")
\textbf{Analytically classified outcomes.} First verify membership in
the maintained domain \(\mathcal D\) and \(B_r(m)>0\) at each node. Pooling existence is

% REPLACES: online_appendix.tex lines 2386-2390 ("the benchmark strength grid the low-cost floor ..." to "... code applies the full rule regardless.")
the benchmark strength grid gross profit at the lowest belief, \(B_r(m)\), is
positive at
every node, entry is positive at every flow, and no price combines
distinct raw posteriors. Constant-posterior atoms remain possible, and
their prices are fixed by the orders, so identity reduces to the
complete order profile there; the code applies the full rule regardless.

% REPLACES: online_appendix.tex lines 2497-2505 ("\textbf{Objects.} Compute m,M ..." to "... economic reversal is absent.")
\textbf{Objects.} Compute \(m,M\), public fundamental bounds
\(\mu_-,\mu_+\), the worst joint buyer posterior \(\phi_-(\mu_-)\), and
the best joint posterior \(\phi_+(\mu_+)\). Compute the four strict
condition margins from \eqref{eq:oa-oa-29}: high-cost
exclusion using the buyer's favorable private signal at weak strength,
high-cost profitability after a favorable strong public price and
private signal, weak residual exclusion, and strong global derivative
margin. Record \(B_{r_1}(\phi_-(\mu_-))\) beside them as
\texttt{low\_cost\_margin}. It measures how strictly the zero-cost realization
prepares at the worst joint posterior; it is not a condition in
\eqref{eq:oa-oa-29}, and its sign carries no theorem label on its own.
If a condition margin fails, withhold the analytical theorem label; do not
conclude that the economic reversal is absent.

% REPLACES: online_appendix.tex lines 2583-2596 ("When the intervals are nonempty choose ..." to "... theorem inequalities explicitly.")
the intervals are nonempty choose

\[
\begin{aligned}
k&=\frac{\Delta_T(r_0)+(1-1/b)\rho m\Delta_T(r_1)}2,\\
c_H&=\frac{B_{r_0}(1/2)+B_{r_1}(M)}2.
\end{aligned}
\tag{OA.78}\label{eq:oa-oa-78}
\]

No low cost is chosen: it is zero in the declaration, and the construction
records \(B_{r_1}(m)\) so that \(0<B_{r_1}(m)<c_H\) can be checked on the row.
Every chosen value is an output of the construction, not a fixed
calibration imposed on all ratios. Check \(0<c_H\), support, and all
theorem inequalities explicitly. Keep the sequence of failed and

% REPLACES: online_appendix.tex lines 2643-2649 ("A threshold scan is first ..." to "... changing other primitives.")
% Lines 2630-2632 in this subsection belong to the noise rewrite and are not
% touched here.
A threshold scan is first an information-experiment comparison. To
interpret a row as an equilibrium with changed preparation cost, set
\(c_H(\tau)=g_L+\tau(g_H-g_L)\), keep the other primitives fixed, and
recheck \(B_r(m)>0\), the ordering \(0<c_H(\tau)\), and the global full-order bound.
Label only validated rows as such. At the strong benchmark the
full-order bound does not depend on the high-cost threshold, but that
fact must not be presumed after changing other primitives.

% REPLACES: online_appendix.tex line 3482 ("numerics/exercises/c6b_price_pools.py")
% Delete this line from the run list. The retired exercise C.6b has no script in
% the release run. The lines for c6_reserve.py and c6c_reserve_events.py stay.

